
@@M0@@

La notion de complexe @@M1@@-parfait envisagée ici est locale, c'est-à-dire qu'il s'agit de perfection relativement à la catégorie fibrée sur @@M2@@ définie par @@M3@@. Cette notion de perfection a bien un sens, car @@M4@@ vérifie les axiomes de (I 4.0) : seule la stabilité par noyau d'épimorphisme n'est pas immédiate ; elle résulte de 1.5.1.

Comme @@M5@@ est quasi-compact, il suffit de vérifier les deux assertions de 2.6 localement sur @@M6@@, de sorte qu'on peut supposer @@M7@@ affine. Le faisceau inversible @@M8@@ est alors ample, et on dispose de résolutions globales sur @@M9@@ (II 2.2.9). Il existe donc un complexe @@M10@@ de @@M11@@, strictement parfait, et un isomorphisme dans @@M12@@ : @@M13@@, ce qui permet de supposer @@M14@@ strictement parfait. Or nous avons vu en 2.4 que, pour @@M15@@ assez grand, les @@M16@@, qui sont tous nuls sauf un nombre fini, sont @@M17@@-acycliques; on a par conséquent~:
@@M18@@
si @@M19@@ est assez grand. Or d'après 1.3 et 1.6, on peut trouver @@M20@@ tel que pour @@M21@@ la condition précédente soit remplie, et que les @@M22@@ soient des @@M23@@-Modules @@M24@@-parfaits pour tout @@M25@@. Le complexe @@M26@@ est alors @@M27@@-parfait, d'où le i).

Si @@M28@@ est choisi tel que précédemment, l'homomorphisme (2.5.3) est alors l'homomorphisme~:
